\documentclass[10pt,a4paper]{amsart}
\usepackage[margin=19mm]{geometry}
\usepackage{amsmath,amssymb,mathtools}
\usepackage[scr=rsfs,cal=euler]{mathalfa}
\usepackage{xfrac}
\usepackage[unicode,hidelinks,bookmarks=false]{hyperref}
\newcommand{\ie}{i.e.\ }
\newcommand{\cf}{cf.\ }
\newcommand{\resp}{respectively\ }
\newcommand{\Subsection}{\subsection}
\providecommand{\nobreakdash}{\nobreak\hbox{-}}
\setlength{\emergencystretch}{2em}
\multlinegap0pt
\usepackage{amssymb}
\usepackage{tikz}
\usepackage{xcolor}
\newtheorem{theorem}{Theorem}
\newtheorem{lemma}{Lemma}
\newcommand{\R}{\mathbb{R}}
\newcommand{\ZR}{\mathbb{R}}
\title{Near diagonal additive energy bound for points on algebraic surfaces}
\author{Yifan Jing, Shukun Wu}
\date{Source: arXiv:2608.14467v2 (CC BY 4.0)}
\begin{document}
\maketitle

\noindent\emph{Source and licence.} Near diagonal additive energy bound for points on algebraic surfaces. Version arXiv:2608.14467v2 (CC BY 4.0), distributed by arXiv under the Creative Commons Attribution 4.0 International licence, http://creativecommons.org/licenses/by/4.0/, as recorded in the arXiv licence metadata for this exact version and shown on the abstract page https://arxiv.org/abs/2608.14467v2. Full licence notice: \texttt{licenses/arxiv-2608.14467-CC-BY-4.0.txt}.\par\smallskip

\noindent\textit{Editorial scope.} Counted unit: the article's lemma energy (source0.tex lines 786--800). The article's complete original proof of it is delivered with it (source0.tex lines 802--892, one \texttt{proof} environment of 91 lines). No mathematics is written, repaired or re-derived: the statement and the proof are the article's own lines, in the article's own notation, and no retained line is substituted or re-flowed.  Retained uncounted as the source-local context the proof needs: lines 610--624.  Nothing was omitted from any retained span, so the package carries no omission record.  No retained line contains a citation, so the wrapper carries no bibliography and no bibliography entry.  No retained line contains a figure environment, an image-inclusion command or drawing code, so no image is delivered and the package declares no asset dependency.  Editorial formatting only, in addition: an AMS \texttt{amsart} preamble that restates the article's own declarations and macros used by the retained lines, namely the environments \texttt{theorem}, \texttt{lemma} on separate counters, together with the standard packages those lines need.  The retained environments print in this document as Theorem 1, Lemma 1 in order of appearance, whereas the article prints the same items with the numbering of its own full document.  The complete native source of the article as distributed in the arXiv e-print for this exact version is bundled unchanged as \texttt{sources/207672/source0.tex}.
\begin{theorem}[B\'ezout]\label{thm: Bezout}
Let $P,Q\in\R[x_1,x_2,x_3]$ be nonconstant and coprime, with degrees $d_1$ and $d_2$, respectively.
There is a possibly empty family of irreducible affine algebraic curves $C_1,\dots,C_r$ defined over $\R$ such that
\begin{equation}
\label{eq:bezout-decomposition}
    Z(P,Q)=\bigcup_{i=1}^r C_i(\R),
\end{equation}
and $\sum_{i=1}^r\deg C_i\leq d_1d_2$.
In particular, if $C_1$ and $C_2$ are distinct irreducible affine algebraic curves defined
over $\R$, then
\begin{equation}
\label{eq:curve-bezout}
    \#\bigl(C_1(\R)\cap C_2(\R)\bigr)  \leq(\deg C_1)(\deg C_2).
\end{equation}
\end{theorem}

\begin{lemma}
\label{lem: wall energy} 
Let $F:\ZR^3\to\ZR$ be a polynomial that is irreducible over $\ZR$ and $\deg_z F\ge 1$. 
Let $q\in\R[u_1,u_2]$ be nonzero of degree $\leq D$,  and set $Q(\vec{u},z)=q(\vec{u})$ so that $Q$ is independent of $z$. 
Let $Y\subset \ZR^3$ and
\[
    W\subset Z(F,Q)
\]
be finite sets with $\#W,\#Y\leq N$. 
Then
\begin{equation}
\label{eq:wall-energy}
    E(W,Y)\leq C_{m_F}\Lambda_{F}(W)D^2N^2.
\end{equation}
\end{lemma}

\begin{proof}
If $q$ is constant, then it is a nonzero constant and $W\subset Z(Q)=\varnothing$, so the result is trivial.
Henceforth assume that $\deg q\ge1$, so $D\geq 1$.
Since $F$ depends nontrivially on $z$ and since $F$ is irreducible over $\R$, the polynomials $F$ and $Q$ are coprime.
Now B\'ezout's Theorem (Theorem~\ref{thm: Bezout}) implies that there are irreducible affine algebraic curves $C_1,\dots,C_r$ defined over $\R$ such that
\[
    Z(F,Q)=\bigcup_{i=1}^rC_i(\R).
\]
Denote by $R_0 = \sum_{i=1}^r\deg C_i$, so 
\[
    R_0\leq m_FD.
\]
Assign every point of $W$ to one curve containing it.  
After discarding empty pieces, this gives a disjoint decomposition
\[
    W=\bigsqcup_{i=1}^r W_i
\]
with $W_i\subseteq C_i(\R)$. 

\smallskip

Let $I$ be the set of indices such that for each $i\in I$, $C_i$ is an affine line. 
Define
\[
    W_{\mathrm{line}}=\bigsqcup_{i\in I}W_i.
\]
Since every such line is contained in $Z(F)$, we have $\#I\leq R_0$, and
\begin{equation}
\label{eq: wall line points}
    \#W_{\mathrm{line}}\leq R_0\Lambda_{ F}(W).
\end{equation}


\smallskip

For $h\in\R^3$, define the translation map $\tau_h(C):=C+h$.  
A pair $(w_i,w_j)\in W_i\times W_j$ with $w_i-w_j=h$ forces
\[
    w_i\in C_i(\R)\cap\tau_h(C_j)(\R).
\]
Thus, for a given $h$ and a pair $(C_i, C_j)$ with $C_i\neq\tau_h(C_j)$, Theorem~\ref{thm: Bezout} implies
\begin{equation}
\label{not-translation}
    \#\{(w_i, w_j)\in W_i\times W_j: w_i-w_j=h \}\leq (\deg C_i)(\deg C_j).
\end{equation} 

Next, consider a pair $(C_i,C_j)$ that is not a pair of affine lines. 
Then there is at most one vector $h$ such that $C_i=\tau_h(C_j)$. 
Indeed, if both $h$ and $h'$ satisfy $C_i=\tau_h(C_j)$ and $C_i=\tau_{h'}(C_j)$, then $C_i$ would be invariant under translation by the nonzero vector $a=h'-h$.
Iterating it gives $C_i = C_i + na$ for every $n\in\mathbb Z$. 
Thus, for any polynomial $P$ that vanishes on $C_i$ and for every $w\in C_i$, the univariate polynomial $g(x)=P(w+xa)$ vanishes on every integer, hence is identically $0$.
Hence $w+\R a\subset C_i(\R)$.
Since an irreducible affine curve containing an affine line must equal that line, $C_i(\ZR)$ is a line.
Similarly, $C_j(\ZR)$ is also a line, and this contradicts the choice of the pair $(C_i, C_j)$. 


Let $T$ be the set of all such exceptional translations $h$ arising from ordered pairs $(C_i,C_j)$ that are not both affine lines and satisfy $C_i=\tau_h(C_j)$.
Then the aforementioned discussion implies
\[
    \#T\leq r^2\leq R_0^2.
\]
Let $d_*(h)$ denote the number of pairs $(w_i,w_j)$ in $W\times W$ such that $w_i-w_j=h$ and the associated component pair $(C_i,C_j)$, with $w_i\in C_i$ and $w_j\in C_j$, is not a pair of two affine lines.
For $h\notin T$, by \eqref{not-translation} and summing over the relevant ordered pairs, 
\[
    d_*(h)\leq \sum_{i,j}(\deg C_i)(\deg C_j)=\Big(\sum_{i=1}^r\deg C_i\Big)^2\leq R_0^2.
\]
For $h\in T$, we trivially have $d_*(h),d_{Y}(h)\leq N$. Since $\sum_hd_{Y}(h)=|Y|^2\leq N^2$, 
\begin{align}
\nonumber
    \sum_h d_*(h)d_{Y}(h)=&\sum_{h\not\in T} d_*(h)d_{Y}(h)+\sum_{h \in T} d_*(h)d_{Y}(h)\\  \label{eq:wall-nonlinear-contribution}
    \leq &\, R_0^2N^2+(\#T)N^2\leq 2R_0^2N^2.
\end{align}

Let $d_{r}(h)$ count the number of remaining pairs, that is, the number of $(w_i,w_j)\in W\times W$ with $h=w_i-w_j$ such that the two assigned components $C_i, C_j$ with $w_i\in C_i, w_j\in C_j$ are both affine lines.
Then
\[
    \sum_h d_{r}(h)=(\#W_{\mathrm{line}})^2.
\]
Since $d_{Y}(h)\leq N$ for each $h$, by \eqref{eq: wall line points}, we obtain
\begin{equation}
\label{eq:wall-line-contribution}
    \sum_h d_{r}(h)d_{Y}(h) \leq N(\#W_{\mathrm{line}})^2
    \leq R_0\Lambda_{ F}(W) N^2.
\end{equation}
Note that
\[
    E(W,Y)=\sum_h\bigl(d_*(h)+d_{r}(h)\bigr)d_{Y}(h),
\]
Combining \eqref{eq:wall-nonlinear-contribution} and
\eqref{eq:wall-line-contribution}, and using $R_0\leq m_FD$, we prove \eqref{eq:wall-energy}.
\end{proof}

\end{document}
